\documentclass[11pt]{article}
\usepackage[T1]{fontenc}
\usepackage{lmodern,amsmath,amssymb,amsthm,mathtools}
\usepackage[margin=1.05in]{geometry}
\usepackage[hidelinks]{hyperref}
\newtheorem{theorem}{Theorem}[section]
\newtheorem{lemma}[theorem]{Lemma}
\newtheorem{proposition}[theorem]{Proposition}
\newtheorem{corollary}[theorem]{Corollary}
\newcommand{\R}{\mathbb R}\newcommand{\tr}{\operatorname{tr}}
\newcommand{\code}[1]{{\small\nolinkurl{#1}}}
\title{Exact calibration tests for strain--pressure identities}
\author{Project Stardust}
\date{Reviewed diagnostic draft, October 29, 2026}
\begin{document}
\maketitle
\begin{abstract}
We assemble three exact diagnostics for velocity gradients.
A family of nonnormal restricted Euler matrices has fixed full spectrum but
variable strain spectrum, with arbitrarily small positive normalized
stretching along trajectories whose raw stretching diverges. Its affine
realization solves the local Navier--Stokes equations with quadratic
pressure, but has infinite whole-space energy. A trace computation then
detects an incompatible mixture of full and trace-free pressure
conventions: the proposed strain right side has trace equal to the squared
Frobenius norm of the velocity gradient under the pressure Poisson
constraint. Finally a smooth periodic viscous ABC flow has zero averaged
stretching derivative, while a displayed plus-pressure balance gives
three at the initial time. Complete elementary proofs reconstruct frozen
formal results. These are calibration examples in classical model
families, not new fluid existence or singularity results; priority for
the diagnostic formulations remains unresolved.
\end{abstract}

\section{Conventions and purpose}
For a real $3\times3$ matrix $A$, use
\[
 S=\tfrac12(A+A^{\mathsf T}),\qquad W=A-S,\qquad
 \omega=(A_{32}-A_{23},A_{13}-A_{31},A_{21}-A_{12}).
\]
For an actual velocity $u$, $A_{ij}=\partial_j u_i$ and $\omega=\nabla\times u$.
The pressure Hessian is $P=\nabla^2p$, so the incompressible pressure
constraint is
\begin{equation}\label{eq:poisson}
 -\tr P=\tr(A^2).
\end{equation}
This convention agrees with the velocity-gradient formulation in
Cantwell~\cite{Cantwell} and the explicit pressure equations of
Yang, Xu, Pumir and He~\cite{Yang}. Full pressure and its deviatoric
part $P_0=P-\tfrac13\tr(P)I$ must be distinguished.

The first test concerns the matrix ODE
\begin{equation}\label{eq:RE}
 A'=-A^2+\tfrac13\tr(A^2)I,\qquad \tr A=0.
\end{equation}
The second is an algebraic test of a specified strain right side. The
third tests a specified spatially averaged identity on an actual smooth
periodic Navier--Stokes flow. Each test has its own hypotheses; neither
an ODE trajectory nor an infinite-energy affine field is promoted to a
finite-energy singularity theorem.

\section{Fixed full spectrum and arbitrarily small normalized stretching}
For $p>2$ set $a=\sqrt{(p-2)(p+1)}>0$ and define
\begin{equation}\label{eq:M}
 M_{p,a}=\begin{pmatrix}-p&0&a\\0&1&0\\-a&0&p-1\end{pmatrix},
 \qquad c(t)=(1-t)^{-1},\qquad A(t)=c(t)M_{p,a}.
\end{equation}
For $\omega\ne0$ let
\[
 \sigma(A)=\frac{\omega^{\mathsf T}S\omega}{|\omega|^2},\qquad
 \rho(A)=\frac{\sigma(A)}{\|S\|_F}.
\]
The normalization uses the Frobenius norm of the symmetric strain,
not the spectral radius of the nonsymmetric full matrix.

\begin{theorem}\label{thm:family}
For every $p>2$ the trajectory \eqref{eq:M} solves \eqref{eq:RE} for
$t<1$. Its strain eigenvalues in increasing order are
$c(t)(-p,1,p-1)$, and its nonzero vorticity is aligned with the
intermediate eigendirection. Its full characteristic polynomial is
\[
 \det(rI-A(t))=(r-c(t))^2(r+2c(t)),
\]
whereas
\begin{equation}\label{eq:ratios}
 \sigma(A(t))=c(t),\qquad
 \rho(A(t))=\frac1{\sqrt{p^2+1+(p-1)^2}}.
\end{equation}
For every $\epsilon>0$ a choice of $p>2$ makes
$0<\rho(A(t))<\epsilon$ for all $0\leq t<1$, while
$\sigma(A(t))\to\infty$ as $t\uparrow1$. No positive uniform lower
threshold for this normalized quantity holds on these exact profiles.
\end{theorem}
\begin{proof}
The middle coordinate of $M$ is one. On the first and third coordinates,
direct multiplication and $a^2=p^2-p-2$ give
\[
 M^2+M=2I,\qquad \tr M=0,\qquad \tr M^2=6.
\]
Thus $-M^2+\tr(M^2)I/3=M$. Since $c'=c^2$ and the ODE right side
is homogeneous of degree two, $A'=c^2M$ is precisely that right side.
Also
\[
 S(M)=\operatorname{diag}(-p,1,p-1),\qquad \omega(M)=(0,2a,0).
\]
Multiplying by $c>0$ proves the strain ordering and alignment and gives
$|\omega(A)|^2=4a^2c^2$ and $S(A)\omega(A)=c\omega(A)$.
The stretching formulas \eqref{eq:ratios} follow, since
$\|S(A)\|_F^2=c^2(p^2+1+(p-1)^2)$.

The determinant factors as
\[
 \det(rI-M)=(r-1)((r+p)(r-p+1)+a^2)=(r-1)^2(r+2).
\]
Rescaling gives the stated polynomial for $A$. The matrix is nonnormal:
the $(1,3)$ entries of $M^{\mathsf T}M$ and $MM^{\mathsf T}$ are
$-a(2p-1)$ and $a(2p-1)$, respectively, and are unequal.
Hence the varying symmetric spectrum does not contradict the fixed
full spectrum.

For $\epsilon>0$, choose $p=2+\epsilon^{-1}$ and the specified $a$.
The denominator in \eqref{eq:ratios} is strictly larger than $p$, and
$\epsilon p=2\epsilon+1>1$, so $0<\rho<1/p<\epsilon$.
Meanwhile $c(t)\to\infty$ as $t\uparrow1$, proving raw stretching
divergence along each selected trajectory.
\end{proof}

For example $p=3,a=2$ gives strain ratios $-3:1:2$; $p=4,a=\sqrt{10}$
gives $-4:1:3$ and normalized stretching $1/\sqrt{26}$. A particular
strain ratio is therefore not universal even within this elementary
self-similar family. Cantwell's general restricted-Euler solution already
distinguishes full-gradient and strain/rotation geometry~\cite{Cantwell}.
The examples here are exact calibration data, not a new solution theory.

\section{An affine realization and its spatial limitation}
\begin{proposition}\label{prop:affine}
With the same parameters, define on $\R^3$, for $t<1$,
\[
 u(t,x)=c(t)M_{p,a}x,\qquad p_{\rm aff}(t,x)=-c(t)^2|x|^2.
\]
These smooth fields satisfy
\[
 \nabla\cdot u=0,\qquad
 \partial_tu+(u\cdot\nabla)u+\nabla p_{\rm aff}=\nu\Delta u
\]
for every fixed $\nu\geq0$, pointwise on this time interval. They are
not periodic and $u(t,\cdot)\notin L^2(\R^3)$.
\end{proposition}
\begin{proof}
The spatial gradient is $cM$, the divergence is $c\tr M=0$, and
$\Delta u=0$. Moreover
\[
 \partial_tu=c^2Mx,\qquad
 (u\cdot\nabla)u=c^2M^2x,\qquad \nabla p_{\rm aff}=-2c^2x.
\]
The momentum identity is therefore $c^2(M+M^2-2I)x=0$.
Its pressure Hessian is $P=-2c^2I$ and satisfies
$-\tr P=6c^2=\tr(A^2)$, so the Poisson convention is consistent.
The middle velocity coordinate is $u_2=c x_2$, which is not periodic.
Since $c>0$, integrating $|u|^2\geq c^2x_2^2$ over expanding boxes
gives an infinite integral. Thus these fields do not meet the
whole-space finite-energy hypothesis.
\end{proof}

They also provide a local pressure-sign test. Define
$q=\omega^{\mathsf T}S\omega$, $b=|S\omega|^2$ and
$h=\omega^{\mathsf T}P\omega$ for this affine field. Directly,
\[
 q=4a^2c^3,\qquad b=4a^2c^4,\qquad h=-8a^2c^4.
\]
Since $q$ is spatially constant and $c'=c^2$,
\[
 (\partial_t+u\cdot\nabla)q=12a^2c^4=b-h\ne b+h.
\]
This equality is a computation for these fields, not a proposed general
formula with all other terms suppressed. It calibrates the sign of the
displayed pressure contraction locally; it supplies no finite global
integrals of the affine field.

\section{A trace obstruction to mixed pressure conventions}
The strain right side to be tested is the explicit matrix expression
\begin{equation}\label{eq:bad}
 R_+(A,P)=-S^2-W^2-P+
 \frac{\tr(S^2)+|\omega|^2/2}{3}I.
\end{equation}
It is this expression, including the isotropic plus sign and the
meaning of $P$, that the next result excludes as a trace-preserving
incompressible strain evolution.

\begin{theorem}\label{thm:trace}
For every real $A,P$ satisfying \eqref{eq:poisson},
\begin{equation}\label{eq:badtrace}
 \tr R_+(A,P)=\|A\|_F^2.
\end{equation}
In particular it is strictly positive if $A\ne0$. Adding a trace-free
matrix correction cannot restore trace zero. The full-pressure
expression and the trace-free-pressure expression
\begin{align}
 R_{
 \rm full}(A,P)&=-S^2-W^2-P,\label{eq:full}\\
 R_{\rm dev}(A,P_0)&=-S^2-W^2-P_0+
 \frac{\tr(S^2)-|\omega|^2/2}{3}I\label{eq:dev}
\end{align}
are trace-consistent under their respective conventions. Under
\eqref{eq:poisson}, \eqref{eq:full} and \eqref{eq:dev} agree.
\end{theorem}
\begin{proof}
The symmetric/skew splitting is Frobenius orthogonal. The skew matrix
associated with the above vorticity has squared Frobenius norm
$|\omega|^2/2$, and hence
\begin{align*}
 \tr W^2&=-|\omega|^2/2,\\
 \tr A^2&=\tr S^2-|\omega|^2/2,\\
 \|A\|_F^2&=\tr S^2+|\omega|^2/2
             =\tr A^2+|\omega|^2.
\end{align*}
For example $|\omega|^2$ is the sum over $i<j$ of
$(A_{ij}-A_{ji})^2$, while the off-diagonal contribution to
$\tr A^2$ is $2\sum_{i<j}A_{ij}A_{ji}$; adding proves the last identity
directly. Taking the trace of \eqref{eq:bad} gives
$|\omega|^2-\tr P$. Substituting \eqref{eq:poisson} proves
\eqref{eq:badtrace}. The Frobenius squared norm is positive unless all
entries vanish, and a trace-free addition leaves this trace unchanged.

The trace of \eqref{eq:full} is $-\tr A^2-\tr P$, which is zero under
\eqref{eq:poisson}. The trace of \eqref{eq:dev} is zero because
$\tr P_0=0$ and $\tr(S^2)-|\omega|^2/2=\tr A^2$.
Finally $P=P_0+\tr(P)I/3=P_0-\tr(A^2)I/3$ under the constraint;
inserting this in \eqref{eq:full} gives \eqref{eq:dev}.
\end{proof}

If $\tr A=0$, the strain trace must remain zero along an incompressible
evolution. Theorem~\ref{thm:trace} thus gives an immediate nonzero-state
obstruction to \eqref{eq:bad}. The theorem is a consistency test, not
a derivation or closure of a nonlocal pressure law.

\section{A periodic positive-viscosity pressure-sign test}
On the $2\pi$-periodic three-torus let
\[
 U(x,y,z)=(\sin z+\cos y,\ \sin x+\cos z,\ \sin y+\cos x).
\]
This is the classical unit-parameter ABC profile; the same convention
is used in~\cite{XYZ}. Write $\langle f\rangle=(2\pi)^{-3}
\int_{[0,2\pi]^3}f$ for normalized spatial average. For a smooth
velocity $v$ with strain $S_v$ and vorticity $\omega_v$, define
\[
 Q(v)=\langle\omega_v^{\mathsf T}S_v\omega_v\rangle,\quad
 B(v)=\langle|S_v\omega_v|^2\rangle,\quad
 H(v,p)=\langle\omega_v^{\mathsf T}(\nabla^2p)\omega_v\rangle,
\]
and the specific viscous contraction
\[
 J(v)=\left\langle\sum_{k,i,j}(S_v)_{ij}
               (\partial_k\omega_{v,i})(\partial_k\omega_{v,j})\right\rangle.
\]
These are ordinary derivatives and integrals of smooth periodic fields.

\begin{theorem}\label{thm:ABC}
For each $\nu>0$, the fields
\[
 u(t)=e^{-\nu t}U,\qquad p(t)=-\tfrac12|u(t)|^2
\]
are a smooth periodic Navier--Stokes solution for all $t\geq0$.
They satisfy $Q(u(t))=0$ for every $t$, whereas at $t=0$,
\begin{equation}\label{eq:ABCvalues}
 B(U)=\tfrac32,\qquad H(U,-|U|^2/2)=\tfrac32,\qquad J(U)=0.
\end{equation}
In particular the identity
\begin{equation}\label{eq:plusbalance}
 \frac{d}{dt}Q(u(t))=B(u(t))+H(u(t),p(t))-2\nu J(u(t))
\end{equation}
fails at $t=0$: its sides are zero and three respectively. Reversing the
pressure sign makes the two sides agree on this test at $t=0$, without
asserting a general evolution formula from this agreement.
\end{theorem}
\begin{proof}
Differentiation gives $\nabla\cdot U=0$, $\nabla\times U=U$ and
$\Delta U=-U$. For a Beltrami field,
$(U\cdot\nabla)U=\nabla K$, $K=|U|^2/2$, by the vector identity
$(U\cdot\nabla)U=\nabla K-U\times(\nabla\times U)$.
Moreover $S_U U=\nabla K$, because the skew part of $\nabla U$
annihilates its vorticity $U$. For $u=e^{-\nu t}U$, the nonlinear
advection is canceled by $\nabla p$, and
$\partial_tu=-\nu u=\nu\Delta u$. This proves the PDE assertion.

Periodic integration by parts and $\nabla\cdot U=0$ give
$Q(U)=\langle U\cdot\nabla K\rangle=0$. At time $t$, stretching is
scaled by $e^{-3\nu t}$, so $Q(u(t))=0$ and its derivative is zero.
The kinetic gradient is
\[
 \nabla K=(\cos x\cos z-\sin x\sin y,
 \ \cos y\cos x-\sin y\sin z,
 \ \cos z\cos y-\sin z\sin x).
\]
Each component has mean square $1/4+1/4=1/2$; the cross terms have
zero average by separation of variables. Since $S_UU=\nabla K$,
this proves $B(U)=3/2$.
For the pressure contraction, integrating $U_iU_j\partial_i\partial_j(-K)$
by parts in $i$ gives
\[
 H(U,-K)=-\langle((U\cdot\nabla)U)\cdot\nabla(-K)\rangle
 =\langle|\nabla K|^2\rangle=\tfrac32.
\]
All boundary terms cancel by periodicity.

It remains to check the precise viscous contraction, rather than omit it.
The strain has zero diagonal and
\[
 S_{12}=\tfrac12(\cos x-\sin y),\quad
 S_{13}=\tfrac12(\cos z-\sin x),\quad
 S_{23}=\tfrac12(\cos y-\sin z).
\]
As $\omega_U=U$, the three terms
$(\partial_kU)^{\mathsf T}S_U(\partial_kU)$, for $k=x,y,z$, are
\begin{align*}
 &-(\cos y-\sin z)\cos x\sin x,\\
 &-(\cos z-\sin x)\sin y\cos y,\\
 &-(\cos x-\sin y)\cos z\sin z.
\end{align*}
Each has zero normalized average; for example integration in $x$ kills
the first. Thus $J(U)=0$. These values prove both the failure of
\eqref{eq:plusbalance} and the limited sign-reversal calibration.
\end{proof}

The affine sign test and periodic test address different possible
loopholes. The former is local and has no finite whole-space energy;
the latter is smooth, periodic, and has arbitrary positive viscosity.
Neither asserts that all pressure contractions have one sign, or that
changing one sign establishes a complete general strain balance.

\section{Formal correspondence and contribution status}
\begingroup\raggedright
Every formal result is frozen at commit
\code{bab688bb6140cd6316c80680095a914ea0d346db}.
The family and its small-normalized-stretching quantifiers are in
\code{RestrictedEulerAlignedGrowth},
\code{RestrictedEulerSpectralDistinction} and
\code{RestrictedEulerArbitraryDepletion}. Actual affine derivatives,
momentum and pressure appear in \code{AffineNavierStokesProfiles} and
\code{AffineStretchingPressureBalance}. The trace test is in
\code{PressureStrainTraceObstruction}. The real periodic derivatives,
means and viscous contraction are assembled in
\code{ABCNavierStokesFlow}, \code{RealBeltramiMeans},
\code{BeltramiViscousTrace} and \code{PeriodicPressureSignCounterexample}.
Exact declaration locations, hashes and human proof reorganizations
are in the accompanying theorem map. In particular the direct
trigonometric proof of $J(U)=0$ replaces the formal invariant-trace
transport proof. Infinite affine energy is the elementary deduction
displayed above; the mapped source explicitly proves spatial
unboundedness and nonperiodicity, not a named integral theorem.

These examples belong to classical restricted-Euler and Beltrami
theory. Cantwell's primary paper was inspected for its model, solution
scope and distinction between full-gradient and strain geometry.
The primary 2024 pressure-Hessian paper fixes the full/deviatoric
conventions; no conflict with its stated equations is alleged.
The ABC profile attribution was checked against the primary
Xin--Yu--Zlato\v{s} model statement. The examples test the explicit
formulas written here, rather than an unspecified published theorem.
Their value is reproducible diagnostic mathematics and a formalization
companion; novelty or priority for any formulation remains unresolved.
No Navier--Stokes regularity, finite-energy blow-up or Millennium
equivalence claim follows. Complete baseline verification evidence was
hash-checked without a redundant full Lean build.
\par\endgroup

\begin{thebibliography}{9}
\bibitem{Cantwell} B. J. Cantwell,
\emph{Exact solution of a restricted Euler equation for the velocity
gradient tensor}, Phys. Fluids A 4 (1992), 782--793.
\href{https://doi.org/10.1063/1.858295}{doi:10.1063/1.858295}.
\bibitem{Yang} P.-F. Yang, H. Xu, A. Pumir and G. W. He,
\emph{Structure and role of the pressure Hessian in regions of strong
vorticity in turbulence}, J. Fluid Mech. 983 (2024), R2.
\href{https://doi.org/10.1017/jfm.2024.143}{doi:10.1017/jfm.2024.143}.
\bibitem{XYZ} J. Xin, Y. Yu and A. Zlato\v{s},
\emph{Periodic orbits of the ABC flow with $A=B=C=1$},
\href{https://arxiv.org/abs/1601.02724}{arXiv:1601.02724} (2016).
Profile comparison only; no orbit result is used here.
\end{thebibliography}
\end{document}
